\chapter{医療統計の問題をどう読むか}\label{ch:01}

\begin{goalbox}
\begin{itemize}
\item 問題文から，アウトカムの型（連続・2 値・カウント・生存時間・反復測定）と，観測単位の独立性・対応を読み取れる．
\item ランダム化比較試験か観察研究かを判別し，交絡の可能性を指摘できる．
\item 推定したい量（効果指標）を式で書き，推定と検定の役割を区別できる．
\item P 値と信頼区間を正しい日本語で解釈できる（「有意差なし」と「差がない」を区別できる）．
\item 群ごとの信頼区間の重なりと2標本検定の結果の関係を導出できる．
\end{itemize}
\end{goalbox}

\begin{exambox}
\begin{itemize}
\item \kakomon{2016}{2}〔2〕〔4〕：$P=0.077$ の結果から言えること，傾向性検定で $P=0.042$ になったときの解釈．
\item \kakomon{2016}{5}（共通）：群ごとの95\%信頼区間の重なりと2標本 $t$ 検定の関係（重なりが区間幅の約29\%未満なら有意）．
\item \kakomon{2022}{1}〔5〕：ハザード比の95\%信頼区間が1を含むときの解釈．
\item \kakomon{2016}{1} と \kakomon{2019}{4}：同じ「2つの測定値の比較」でも，対応あり（治療前後）と独立（2群）で手法が変わる．
\end{itemize}
設問に「結果を解釈せよ」「この結果から言えることを述べよ」とあれば，本章の型で答える．
\end{exambox}

\section{問題設定：医療統計の設問の読み方}\label{sec:01-read}
医薬生物学の設問は，ほぼ例外なく次の情報を問題文の中に含んでいる．
答案を書き始める前に，これらを明示的に書き出す習慣をつけるとよい．
\begin{enumerate}[label=(\alph*)]
\item \textbf{研究デザイン}：ランダム化比較試験（並行群間・クロスオーバー），コホート研究，症例対照研究，マッチング，単群試験など．
\item \textbf{観測単位とその関係}：患者，ペア，層．同一患者の反復測定か，独立な2群か．
\item \textbf{アウトカムの型}：連続，2 値，カウント（人時間つき），生存時間（打ち切りつき），反復測定．
\item \textbf{比較の目的}：優越性・同等性・非劣性，関連の有無，予測，用量反応．
\item \textbf{推定対象}：平均差，リスク差，リスク比，オッズ比，ハザード比，AUC，感度など．
\item \textbf{確率モデル}：二項，多項，ポアソン，正規，指数，比例ハザードなど．問題文に明示されることが多い．
\end{enumerate}
(a)〜(c)が決まると候補となる手法はかなり絞られ，(d)〜(f)で一意に決まる．

\section{基本概念}\label{sec:01-basic}
\subsection{アウトカムの種類}\label{subsec:01-outcome}
\begin{table}[htbp]
\centering\small
\caption{アウトカムの型と代表的な要約・手法}\label{tab:01-outcome}
\begin{tabular}{@{}lL{35mm}L{35mm}L{45mm}@{}}
\toprule
型 & 例 & 要約・効果指標 & 代表的な手法（本書の章）\\
\midrule
連続 & 血圧，検査値，AUC（薬物動態） & 平均差，中央値の差，幾何平均比 & $t$ 検定，順位検定（\ref{ch:03}），同等性（\ref{ch:13}）\\
2 値 & 有効／無効，発症の有無 & リスク差，リスク比，オッズ比 & $\chi^2$ 検定，McNemar，層別の検定（\ref{ch:04}），ロジスティック（\ref{ch:06}）\\
カウント & 有害事象件数，死亡数 & 率（人時間あたり），率比 & ポアソンモデル（\ref{ch:17}）\\
生存時間 & 死亡・再発までの時間 & 生存関数，中央生存時間，ハザード比 & KM，ログランク（\ref{ch:09}），Cox・競合リスクの解釈（\ref{ch:10}）\\
反復測定 & 同一患者の複数時点の値 & 時点別平均，変化量 & 混合モデル（\ref{ch:08}）\\
\bottomrule
\end{tabular}
\end{table}

生存時間は「時間という連続量」だが，観察打ち切り（censoring）があるため連続データの手法をそのまま適用できない．
これが生存時間解析という独立した分野が必要になる理由である（\cref{ch:09}）．
カウントは，観察期間（人時間）が個体や集団ごとに異なるので「件数」ではなく「率」を比較する．

\subsection{独立なデータと対応のあるデータ}\label{subsec:01-paired}
\Hissu 同じ患者から得た2つの値（治療前後，左右の眼，2種類の検査），マッチさせたペア，
クロスオーバー試験の2期の値は\emphx{対応のあるデータ}である．
対応があると2つの値は一般に正の相関をもつ．

\begin{meidai}[対応のある差の分散]\label[meidai]{prop:01-pairvar}
$(X_i,Y_i)$ $(i=1,\dots,n)$ が独立に，$\V[X_i]=\sigma_X^2$，$\V[Y_i]=\sigma_Y^2$，$\Corr(X_i,Y_i)=\rho$ に従うとする．
$D_i=Y_i-X_i$ とすると
\begin{equation}
\V[\bar D]=\frac{\sigma_X^2+\sigma_Y^2-2\rho\sigma_X\sigma_Y}{n}.\label{eq:01-pairvar}
\end{equation}
一方，大きさ $n$ の独立な2群の平均の差 $\bar Y-\bar X$ の分散は $(\sigma_X^2+\sigma_Y^2)/n$ である．
\end{meidai}
\begin{proof}
$\V[D_i]=\V[Y_i]+\V[X_i]-2\Cov(X_i,Y_i)$ と $D_i$ の独立性から直ちに従う．
\end{proof}
$\rho>0$ なら対応を利用した比較の方が分散が小さい．
逆に，対応のあるデータに独立2群の手法を使うと，分散を過大評価して検出力を失い，
独立なデータに対応のある手法は使えない（ペアが定義できない）．
2 値データでも同じで，対応があればMcNemar型（\cref{sec:04-mcnemar}），なければ2標本の比率の比較になる．

\subsection{ランダム化比較試験と観察研究}\label{subsec:01-rct}
\emphx{ランダム化}（無作為割付）は，治療の割付を患者の特性と確率的に独立にする操作である．
これにより，測定されたか否かにかかわらず，あらゆる予後因子の分布が群間で（期待値の意味で）等しくなり，
群間差を治療効果として解釈できる．
観察研究では治療の選択が患者の特性に依存するため，群間差に\emphx{交絡}が混入しうる（\cref{ch:07}）．
設問が観察研究なら「交絡の可能性」「調整の方法」に触れることが，解釈の問いでほぼ必ず求められる．

\subsection{推定したい量}\label{subsec:01-estimand}
推定したい母集団の量を先に言葉と式で定めることが，手法選択の出発点である．
たとえば2群の有効率 $\pi_1,\pi_0$ があるとき，
\[
\RD=\pi_1-\pi_0,\qquad \RR=\frac{\pi_1}{\pi_0},\qquad \OR=\frac{\pi_1/(1-\pi_1)}{\pi_0/(1-\pi_0)}
\]
はすべて「差がない」ことを同じ帰無仮説 $\pi_1=\pi_0$ で表すが，推定量・分散・解釈は異なる（\cref{ch:02}）．
臨床試験では，推定対象を母集団・治療・変数・中間事象の扱い・要約指標の5要素で規定する
「estimand」の枠組みが国際的に採用されている（\cref{sec:02-estimand}）．

\section{推定と検定}\label{sec:01-inference}
\subsection{推定と検定の役割}
\begin{itemize}
\item \textbf{点推定と区間推定}：効果の\emph{大きさ}とその不確かさを示す．臨床的判断の主役．
\item \textbf{検定}：帰無仮説（通常「効果なし」）がデータと両立するかを判定する．効果の大きさは示さない．
\end{itemize}
医学研究の報告では，効果指標の点推定値と95\%信頼区間を示し，P 値は補助的に示すのが標準である．

\subsection{P 値}\label{subsec:01-pvalue}
\begin{teigi}[P 値]\label{def:01-pvalue}
検定統計量 $T$（大きいほど帰無仮説から離れる）と観測値 $t_{\mathrm{obs}}$ に対し，
\[
p=\sup_{\theta\in\Theta_0}\Prob_\theta\bigl(T\ge t_{\mathrm{obs}}\bigr)
\]
をP 値という．「帰無仮説が正しいとしたとき，観測された結果と同じかそれ以上に極端な結果が得られる確率」である．
\end{teigi}
\Youchui P 値は「帰無仮説が正しい確率」ではなく，「効果の大きさ」でもない．
同じ効果の大きさでも標本サイズが大きければP 値は小さくなる．

\begin{pointbox}[P 値の解釈の型（有意水準5\%）]
\begin{itemize}
\item $p<0.05$：「有意水準5\%で帰無仮説は棄却され，〇〇に差があると言える．」
      さらに効果の方向と大きさ（点推定値・信頼区間）を添える．
\item $p\ge 0.05$：「有意水準5\%で帰無仮説は棄却されず，〇〇に差があるとは言えない．」
      \textbf{「差がない」「同等である」とは言わない}．検出力不足の可能性にも触れる．
\end{itemize}
\end{pointbox}

\kakomon{2016}{2} では，4群の有効率の一様性の $\chi^2$ 検定で $p=0.077$ となった．
正しい解釈は「すべての群で母有効率が等しいという帰無仮説は棄却されず，
4群間で有効率に差があるとは言えない」である．
同じデータで用量の順序を利用したCochran--Armitage傾向性検定（\cref{sec:04-trend}で扱う）は $p=0.042$ となり有意であった．
対立仮説を「用量とともに単調に変化」という方向に絞ると検出力が上がるためであり，
「どちらの検定を使うか」はデータを見る前に研究目的から決めておくべきことである．

\subsection{信頼区間}\label{subsec:01-ci}
\begin{teigi}[信頼区間]\label{def:01-ci}
データの関数 $L(\bm X),U(\bm X)$ が，すべての $\theta$ で $\Prob_\theta(L\le\theta\le U)\ge 1-\alpha$ を満たすとき，
$[L,U]$ を $\theta$ の信頼係数 $100(1-\alpha)\%$ の信頼区間という．
\end{teigi}
確率的なのは区間 $[L,U]$ の方で，$\theta$ は定数である．
「真の値が95\%の確率でこの区間に入る」は頻度論の解釈としては不正確で，
「同じ手順を繰り返すと，構成される区間の95\%が真の値を含む」が正しい．

\begin{meidai}[検定と信頼区間の双対性]\label[meidai]{prop:01-duality}
各 $\theta_0$ について有意水準 $\alpha$ の検定の受容域を $A(\theta_0)$ とすると，
$C(\bm x)=\{\theta_0:\bm x\in A(\theta_0)\}$ は信頼係数 $1-\alpha$ の信頼区間（信頼集合）である．
逆に信頼区間から検定を作ることもできる．
\end{meidai}
\begin{proof}
$\theta_0\in C(\bm X)\iff \bm X\in A(\theta_0)$ であり，$\Prob_{\theta_0}(\bm X\in A(\theta_0))\ge1-\alpha$ である．
\end{proof}
したがって，同じ統計量（例：Wald 統計量）から検定と区間を作る場合，
「$H_0:\theta=\theta_0$ の両側 $\alpha$ 検定が棄却されない」$\iff$「$100(1-\alpha)\%$ 信頼区間が $\theta_0$ を含む」．
Pearson $\chi^2$ 検定（スコア型）と対数オッズ比の信頼区間（Wald 型．いずれも\cref{ch:04}で扱う）のように作り方が異なると，境界付近で結論が一致しないことがある．
比の指標（RR，OR，HR）では $\theta_0=1$，差の指標（RD，平均差）では $\theta_0=0$ が「効果なし」である．
\Youchui 生物学的同等性の判定には95\%ではなく90\%信頼区間を使う（\cref{ch:13}で扱う）．

\kakomon{2022}{1} では，Cox回帰（\cref{ch:10}で扱う）で $\hat\beta=0.06$，$\SE=0.18$ から
$\HR=e^{0.06}\approx1.06$，95\%信頼区間 $(e^{0.06-1.96\times0.18},e^{0.06+1.96\times0.18})\approx(0.75,1.51)$ を得る．
区間が1を含むので「試験群とプラセボ群で死亡までの時間に差があるとは言えない」と解釈する．
このとき区間の幅も解釈に含めるとよい：$0.75$ から $1.51$ まで，すなわち25\%のハザード減少から51\%の増加まで，
臨床的に重要な差を否定できない程度にデータが少ない，という情報を区間は伝えている．

\subsection{片側と両側}\label{subsec:01-sided}
優越性を示す確認的試験では，片側2.5\%の検定と両側5\%の検定は，棄却限界（$z_{0.025}=1.96$）が一致する．
非劣性・同等性や群逐次デザインでは片側で考える方が自然なことが多い．
設問が片側か両側かを明示している場合はそれに従い，P 値を計算するときは
「両側P 値＝片側P 値×2」（対称な分布の場合）であることに注意する．

\section{数理：信頼区間の重なりと2標本検定}\label{sec:01-overlap}
\Hinshutsu 論文の図でよく見る「群ごとの平均と95\%信頼区間」は，群間比較の検定と同じではない．
\kakomon{2016}{5} の内容を一般化して導く．

2群の大きさがともに $n$，標本標準偏差がともに $s$ とし，$\bar x>\bar y$ とする．
各群の母平均の95\%信頼区間は $\bar x\pm c_1 s/\sqrt n$，$\bar y\pm c_1 s/\sqrt n$，$c_1=t_{n-1}(0.025)$ である．
等分散の2標本 $t$ 検定（プールした分散は $s^2$）は
\begin{equation}
\bar x-\bar y>c_2\sqrt{\tfrac{2}{n}}\,s,\qquad c_2=t_{2n-2}(0.025)\label{eq:01-ttest}
\end{equation}
のとき有意である．2つの区間の重なりの長さを $D$ とすると
\[
D=\Bigl(\bar y+c_1\frac{s}{\sqrt n}\Bigr)-\Bigl(\bar x-c_1\frac{s}{\sqrt n}\Bigr)=2c_1\frac{s}{\sqrt n}-(\bar x-\bar y).
\]
\begin{meidai}\label[meidai]{prop:01-overlap}
(1) 2つの信頼区間が重ならなければ（$D<0$），2標本 $t$ 検定は有意である．\\
(2) 重なりがあっても，
\begin{equation}
\frac{D}{2c_1 s/\sqrt n}<1-\frac{1}{\sqrt2}\cdot\frac{c_2}{c_1}\label{eq:01-overlap}
\end{equation}
すなわち重なりが1本の区間の幅に対してこの割合未満なら有意である．右辺は $n\to\infty$ で $1-1/\sqrt2\approx0.293$ に近づく．
\end{meidai}
\begin{proof}
$t$ 分布の上側点は自由度について減少するので $c_2<c_1$，よって $2c_1>\sqrt2 c_2$．
(1) $D<0$ なら $\bar x-\bar y>2c_1 s/\sqrt n>\sqrt2 c_2 s/\sqrt n$ で\cref{eq:01-ttest}が成り立つ．
(2) \cref{eq:01-ttest}は $D=2c_1s/\sqrt n-(\bar x-\bar y)<(2c_1-\sqrt2c_2)s/\sqrt n$ と同値であり，両辺を $2c_1s/\sqrt n$ で割ればよい．
\end{proof}
「区間が重なっているから有意差はない」という判断は誤りである．
差の検定の標準誤差は $\sqrt2\,s/\sqrt n$ であって，各群の標準誤差の和 $2s/\sqrt n$ ではないからである．

\section{仮定}\label{sec:01-assump}
すべての推測は仮定の上に立つ．記述式では「この検定に通常置かれる仮定を示せ」が頻出である
（\kakomon{2016}{1}，\kakomon{2019}{4}）．仮定は次の3層で整理すると漏れにくい．
\begin{enumerate}
\item \textbf{独立性}：観測単位どうしが独立（対応があるならペア単位で独立）．
\item \textbf{分布の仮定}：正規性，等分散，比例ハザード，対称性など．
\item \textbf{デザイン上の仮定}：ランダム化，無情報な打ち切り（\cref{ch:09}で扱う），欠測のメカニズム（\cref{ch:16}で扱う），交絡の不在．
\end{enumerate}
漸近的手法ではさらに「標本サイズが十分大きい（各セルの期待度数が十分）」が加わる．

\section{結果の解釈}\label{sec:01-interp}
解釈の答案は次の3要素を含めると採点者に伝わりやすい．
\begin{enumerate}
\item \textbf{統計的結論}：有意か否か，信頼区間が帰無値を含むか．
\item \textbf{効果の方向と大きさ}：点推定値と信頼区間を医学の言葉で（例：「死亡のハザードが約30\%低い」）．
\item \textbf{限界}：観察研究なら交絡，小標本なら検出力，多重性，欠測，外的妥当性．
\end{enumerate}
\begin{center}\small
\begin{tabular}{@{}L{45mm}L{100mm}@{}}
\toprule
結果 & 解釈の例\\
\midrule
$\OR=2.0$（95\%CI 1.3--3.1） & 曝露群の発症オッズは非曝露群の約2倍で，区間は1を含まないので有意な関連がある．発症が稀ならリスク比もおよそ2と解釈できる．\\
$\HR=0.70$（95\%CI 0.55--0.89） & 試験群の死亡ハザードは対照群の0.70倍（30\%低下）．比例ハザードが成り立つなら，どの時点でもこの比が保たれる．\\
$\RD=0.05$（95\%CI $-0.02$--$0.12$） & 有効率の差の点推定は5ポイントだが区間が0を含み，差があるとは言えない．$-2$ポイントから12ポイントまでの差はデータと両立する．\\
$p=0.077$ & 5\%水準で有意ではない．「差がない」とは結論できない．\\
\bottomrule
\end{tabular}
\end{center}

\section{手法選択の基本}\label{sec:01-choice}
\begin{pointbox}[手法選択の質問リスト]
\begin{enumerate}
\item アウトカムは何型か（\cref{tab:01-outcome}）．
\item 比較する群は独立か，対応があるか（同一患者・マッチ・クロスオーバー）．
\item 群の数は2か，3以上か．群に順序（用量）があるか．
\item 調整すべき共変量・層はあるか（観察研究か，層別ランダム化か）．
\item 打ち切り・欠測・競合事象はあるか．
\item 標本サイズは大きいか（漸近近似でよいか，exactが必要か）．
\item 目的は優越性・同等性・非劣性・予測のどれか．
\end{enumerate}
\end{pointbox}
これらの答えから手法を選ぶ一覧表を\cref{app:A}に掲げた．

\section{よくある誤り}\label{sec:01-err}
\begin{warnbox}
\begin{itemize}
\item 「$p>0.05$ だから両群は同等」→ 誤り．同等性は，あらかじめ定めた同等域に信頼区間が含まれることで示す（\cref{ch:13}）．
\item 「P 値が小さいほど効果が大きい」→ 誤り．P 値は効果の大きさと標本サイズの両方で決まる．
\item 群ごとの信頼区間が重なることを根拠に「有意差なし」とする → 誤り（\cref{prop:01-overlap}）．
\item 対応のあるデータに2標本 $t$ 検定を使う，またはその逆．
\item 比の指標の信頼区間を，比のスケールで $\hat\theta\pm1.96\SE$ として作る → 誤り．対数スケールで作って指数変換する．
\item データを見てから片側・両側を選ぶ，複数の検定から有意なものだけを報告する．
\end{itemize}
\end{warnbox}

\section{数理統計との接続}\label{sec:01-link}
\begin{linkbox}
\begin{itemize}
\item 本書の漸近的検定のほとんどは，Wald検定・スコア検定・尤度比検定の3つのいずれかである
      （『現代数理統計学の基礎』7.3節）．分母の分散を帰無仮説の下で推定するのがスコア型
      （例：Pearson $\chi^2$（\cref{subsec:04-pearson}），McNemar（\cref{sec:04-mcnemar}），ログランク（\cref{ch:09}）），
      推定値で評価するのがWald型（例：対数オッズ比の信頼区間，\cref{subsec:04-var}）．
\item 比の指標の信頼区間は，対数変換とデルタ法（同書 定理5.20）で作る．
      多変量デルタ法は同書の範囲外なので\cref{sec:04-delta}で導出する．
\item 検定と信頼区間の双対性は，同書では明示されていないが，定義から直ちに導ける（\cref{prop:01-duality}）．
\item P 値の一様性 $p(\bm X)\sim\Unif(0,1)$（連続分布の場合）は確率積分変換による（同書 命題2.19，定理7.22）．
      多重性の議論（\cref{ch:14}）はこの性質を使う．
\end{itemize}
\end{linkbox}

\section{例題}\label{sec:01-ex}
\begin{reidai}[信頼区間の重なり]\label{reidai:01-2}
2群（各 $n=10$）で連続量を測定し，両群の標本標準偏差はともに $s=5.0$ であった．
$t_9(0.025)=2.262$，$t_{18}(0.025)=2.101$ を用いよ．
\begin{enumerate}[label=(\arabic*)]
\item 各群の平均の95\%信頼区間の半幅を求めよ．
\item 2標本 $t$ 検定（両側5\%）が有意になる平均の差の最小値を求めよ．
\item 平均の差が $6.0$ のとき，2つの信頼区間は重なるか．また検定は有意か．
\item 検定が有意になるための，重なりの長さの上限と，それが区間幅に占める割合を求めよ．
\end{enumerate}
\end{reidai}

\begin{reidai}[対応の利用]\label{reidai:01-3}
治療前後の値 $(X,Y)$ は
$\V[X]=\V[Y]=\sigma^2=100$，相関 $\rho=0.8$ をもつとする．
\begin{enumerate}[label=(\arabic*)]
\item 20名の前後差の平均 $\bar D$ の標準誤差を求めよ．
\item 同じ母集団から独立に20名ずつ2群を取り，平均の差をとった場合の標準誤差を求めよ．
\item (1)の標準誤差が(2)より小さくなる $\rho$ の条件を求めよ．
\end{enumerate}
\end{reidai}

\begin{reidai}[解釈]\label{reidai:01-4}
新薬とプラセボの無作為化試験で，6か月以内の再入院割合を比較した．
リスク比の点推定値は $0.80$，95\%信頼区間は $(0.62,1.03)$，両側 $p=0.08$ であった．
区間は対数スケールで $\log\widehat{\RR}\pm1.96\,\SE$ として作られたものとし，
$\log0.80=-0.2231$，$\log0.62=-0.4780$，$\log1.03=0.0296$ を用いてよい．
\begin{enumerate}[label=(\arabic*)]
\item この結果を，統計的結論・効果の方向と大きさ・限界の3要素を含めて解釈せよ．
\item 「新薬は再入院を減らさないことが示された」「帰無仮説が正しい確率は8\%である」という2つの記述の誤りをそれぞれ説明せよ．
\item 対数リスク比の標準誤差 $\SE$ を信頼区間から求めよ．
\item 点推定値が同じで，$\SE$ が症例数の平方根に反比例すると仮定する．95\%信頼区間の上限が1を下回るには，症例数をおよそ何倍にすればよいか．
\end{enumerate}
\end{reidai}

\section{解答}\label{sec:01-sol}
\enlargethispage{1\baselineskip}
\begin{kaitou}{reidai:01-2}
(1) $t_9(0.025)\,s/\sqrt{10}=2.262\times5.0/3.162=3.577$．\\
(2) 差の標準誤差は $s\sqrt{2/10}=5.0\times0.4472=2.236$．有意になるのは差が $t_{18}(0.025)\times2.236=2.101\times2.236=4.698$ を超えるとき．\\
(3) 重なり $D=2\times3.577-6.0=1.154>0$ で区間は重なる．しかし $6.0>4.698$ なので検定は有意である．\\
(4) \cref{prop:01-overlap}より上限は $2\times3.577-4.698=2.456$．区間幅 $7.154$ に対する割合は
$2.456/7.154=0.343$．公式\cref{eq:01-overlap}でも $1-0.7071\times2.101/2.262=1-0.657=0.343$ となる．
自由度が小さいため漸近値 $0.293$ より大きい．
\end{kaitou}

\begin{kaitou}{reidai:01-3}
(1) $\V[D_i]=2\sigma^2(1-\rho)=2\times100\times0.2=40$ より $\SE=\sqrt{40/20}=\sqrt2=1.414$．\\
(2) $\SE=\sqrt{100/20+100/20}=\sqrt{10}=3.162$．\\
(3) (1)の分散 $2\sigma^2(1-\rho)/n$ と(2)の $2\sigma^2/n$ を比べ，$\rho>0$ のとき(1)が小さい．
ただし(1)は前後比較であり，対照群がなければ時間の経過による変化と治療効果を区別できない点に注意する．
\end{kaitou}

\begin{kaitou}{reidai:01-4}
(1) 統計的結論：区間が1を含み $p=0.08\ge0.05$ なので，有意水準5\%で再入院割合に差があるとは言えない．
効果の方向と大きさ：新薬群の再入院リスクはプラセボ群の0.80倍（20\%低い）と推定され，
38\%の低下（0.62倍）から3\%の増加（1.03倍）までの値がデータと両立する．
限界：区間が広く，臨床的に意味のある低下を否定できない．検出力不足の可能性がある．\\
(2) 前者：有意でないことは効果がないことの証明ではない．区間は38\%の低下も含む．
後者：P 値は「帰無仮説が正しいとしたとき，観測値以上に極端な結果が得られる確率」であり（\cref{def:01-pvalue}），帰無仮説が正しい確率ではない．\\
(3) 区間の幅は対数スケールで $0.0296-(-0.4780)=0.5076=2\times1.96\,\SE$ なので $\SE=0.5076/3.92=0.1295$．
（区間の中点 $-0.2242$ は $\log0.80=-0.2231$ とほぼ一致する．また $z=0.2231/0.1295=1.72$ は $p\approx0.08$ と整合する．）\\
(4) 症例数を $k$ 倍にすると $\SE$ は $0.1295/\sqrt k$ になる．上限が1未満 $\iff$ $-0.2231+1.96\times0.1295/\sqrt k<0$
$\iff$ $\sqrt k>0.2538/0.2231=1.138$ $\iff$ $k>1.29$．およそ1.3倍（3割増）の症例数が必要である．
\end{kaitou}

\begin{summarybox}
\begin{itemize}
\item 設問を読んだら，デザイン・観測単位・アウトカムの型・目的・推定対象・確率モデルを書き出す．
\item 対応のあるデータは差（または不一致ペア）で分析する．$\V[\bar D]=(\sigma_X^2+\sigma_Y^2-2\rho\sigma_X\sigma_Y)/n$．
\item P 値は効果の大きさではない．有意でないことは「差がない」ことを意味しない．
\item 信頼区間と検定は双対：比の指標は1，差の指標は0を含むかで判定．比の指標は対数スケールで区間を作る．
\item 群ごとの95\%信頼区間が重なっていても，重なりが幅の約29\%（小標本ではもう少し大きい）未満なら2標本検定は有意．
\item 解釈の答案は「統計的結論＋効果の方向と大きさ＋限界」の3要素で書く．
\end{itemize}
\end{summarybox}
